\chapter{Basic Controls}
\label{ch:basic-controls}

This chapter describes the main operator controls available on the \textit{RC Screen}. The interface is organized into left and right controller sides plus a central telemetry area, allowing you to drive the robot while monitoring system status in real time.

\section{RC Screen Layout}
\label{sec:rc-layout}

The RC Screen is composed of three main regions:

\begin{itemize}
    \item \textbf{Left Controller Side:} primary controls (joystick mode dependent), switches, knob, and a top status widget.
    \item \textbf{Center Display:} a telemetry plot used to visualize live data.
    \item \textbf{Right Controller Side:} secondary controls (joystick mode dependent), switches, knob, and a top status widget.
\end{itemize}

A background panel provides contrast for readability, and the screen is intended to be used in \textbf{landscape orientation} for ergonomic spacing.

\section{Joystick Control}
\label{sec:joystick-control}

Each controller side includes a virtual joystick. The joystick behavior depends on the selected \textbf{stick mode} (configured in user settings):

\begin{itemize}
    \item \textbf{Spring mode:} the joystick returns to center when released (useful for momentary input such as steering or speed trim).
    \item \textbf{Hold mode:} the joystick keeps its position when released (useful for sustained throttle or setpoints).
    \item \textbf{Directional modes:} some configurations may emphasize horizontal vs.\ vertical control depending on the robot's requirements.
\end{itemize}

Typical joystick usage includes:

\begin{itemize}
    \item Forward/backward movement (throttle)
    \item Left/right turning (steering)
    \item Camera pan/tilt (when mapped to a gimbal)
    \item Any continuous control channel required by the robot (e.g., setpoint control)
\end{itemize}

\section{Switches (Discrete Channels)}
\label{sec:switches}

Both left and right sides provide \textbf{toggle switches} for discrete control channels. Switches are intended for actions that are either \textit{on/off} or select between a small number of states, such as:

\begin{itemize}
    \item Enabling/disabling motors
    \item Selecting drive modes (e.g., slow/normal/boost)
    \item Turning auxiliary devices on/off (lights, pump, fan, etc.)
\end{itemize}

\section{Knobs (Analog Channels)}
\label{sec:knobs}

Each side also includes a \textbf{rotary knob} for analog adjustment. Knobs are useful for fine-tuning parameters while driving:

\begin{itemize}
    \item Steering sensitivity or trim
    \item Speed limit
    \item Camera zoom or focus (if supported)
    \item PID gains or other tuning parameters (advanced setups)
\end{itemize}

\section{Center Buttons}
\label{sec:center-buttons}

The RC Screen includes \textbf{push buttons} near the center area (top and bottom per side). These are intended for quick actions that should be triggered immediately, for example:

\begin{itemize}
    \item Horn / buzzer
    \item Reset / acknowledge
    \item Take picture / start recording (camera payload)
    \item Trigger an automated routine
\end{itemize}

\section{Indicators and Telemetry}
\label{sec:telemetry}

The interface provides several forms of feedback so the operator can monitor the system while controlling it:

\begin{itemize}
    \item \textbf{Analog Indicator (top-left widget):} displays a numeric value with a title (e.g., RPM, temperature, signal strength).
    \item \textbf{Battery Status (top-right widget):} shows battery level with a title label.
    \item \textbf{LED Indicators:} a set of LEDs for quick status cues (e.g., armed state, warnings, link status).
    \item \textbf{Panel Indicators:} per-side panels showing a value, an on/off state, a color, and a title (useful for named telemetry channels).
    \item \textbf{Center Plot:} a live plot showing one or more telemetry series over time.
\end{itemize}

These indicators help confirm command reception and allow the operator to detect issues such as low battery, loss of link, or sensor warnings early.

\section{Summary}
\label{sec:basic-controls-summary}

In general:
\begin{itemize}
    \item Use \textbf{joysticks} for continuous motion control.
    \item Use \textbf{switches and buttons} for discrete actions and mode changes.
    \item Use \textbf{knobs} for fine-grained tuning.
    \item Use \textbf{telemetry widgets and plots} to monitor system state while operating.
\end{itemize}